% ============================================================
% CAPITOLO 5 -- VALIDAZIONE QUALITÀ DEI TEST
% ============================================================
\section{Validazione della qualità dei test}\label{sec:adeguatezza}

Le sezioni precedenti hanno applicato quattro famiglie di test alle
due classi selezionate, senza tuttavia confrontarne sistematicamente
l'efficacia. Questo capitolo affronta tale confronto su tre livelli
complementari: quanto codice viene effettivamente eseguito da ciascuna
famiglia di test (metriche di adeguatezza); quanto i test sono in
grado di rilevare difetti introdotti artificialmente nel codice
(mutation testing); e quale stima di affidabilità si può derivare
dall'insieme finale di test disponibile (reliability).


\subsection{Metriche di adeguatezza}

Per confrontare quantitativamente l'efficacia delle quattro famiglie
di test applicate a ciascuna classe (test manuali, Randoop, LLM,
EvoSuite), sono state scelte due metriche di adeguatezza control-flow:
\emph{Statement Coverage} e \emph{Branch Coverage}. La scelta di
queste due metriche consente un confronto omogeneo con i dati già
raccolti per EvoSuite nel capitolo precedente. Per ciascuna
combinazione famiglia/classe, la copertura è stata misurata isolando
l'esecuzione della sola famiglia in questione con JaCoCo, azzerando i
dati di copertura precedenti prima di ogni misurazione. La
\tabref{tab:adeguatezza-confronto} riporta i risultati completi.

EvoSuite risulta la famiglia più efficace su entrambe le classi
(88.5\% e 95.2\% statement coverage rispettivamente), risultato atteso
poiché è l'unica delle quattro tecniche esplicitamente guidata dal
criterio di copertura durante la generazione. La suite manuale (BB) è
invece la meno efficace in termini di copertura strutturale (42.6\% e
66.7\%): questo non indica una debolezza del category-partition come
tecnica, ma riflette il fatto che le categorie individuate mirano a
rappresentare gli scenari d'uso funzionali della classe, non a
massimizzare la copertura di ogni riga o diramazione del codice.

Il dato più rilevante emerso dal confronto è il comportamento di
Randoop, che presenta un divario estremo tra le due classi: 79.0\%
statement coverage su \texttt{Utf8}, contro appena l'1.6\% su
\texttt{Resolver}. La causa è che Randoop costruisce gli input tramite
reflection, riuscendo a generare valori primitivi (\texttt{String},
\texttt{byte[]}) sufficienti a esercitare efficacemente \texttt{Utf8},
ma non riuscendo a costruire per via casuale oggetti \texttt{Schema} e
\texttt{GenericData} semanticamente validi, necessari per attraversare
la logica di \texttt{Resolver.resolve()}. Questo risultato è coerente
con un limite noto del random testing basato su reflection, quando
applicato a domini con oggetti compositi articolati.

Gli LLM si posizionano stabilmente al secondo posto su entrambe le
classi (75.4\%/68.4\% su \texttt{Resolver}, 97.1\%/95.5\% su
\texttt{Utf8}), superando sia la suite manuale sia Randoop: avendo
accesso al codice sorgente e a vincoli di progetto espliciti, la
generazione tramite LLM produce test sistematici che si avvicinano
alla qualità di un approccio guidato dalla copertura, pur senza
esserlo esplicitamente.


\subsection{Mutation testing con PIT}\label{sec:mutation}

Il mutation testing consiste nell'introdurre piccole modifiche
artificiali nel bytecode compilato (\emph{mutanti}) e verificare se i
test esistenti le rilevano: un mutante è detto \emph{ucciso} se
almeno un test fallisce in sua presenza, \emph{sopravvissuto}
altrimenti. A differenza delle metriche di adeguatezza discusse in
precedenza, questa tecnica misura non solo quanto codice viene
eseguito dai test, ma quanto efficacemente le loro asserzioni sono in
grado di distinguere un comportamento corretto da uno alterato.
L'analisi è stata condotta con PIT sulla sola suite di test manuali
(BB), come richiesto dal punto 4b della consegna.

\paragraph{Problemi di integrazione risolti} Sono emersi due problemi
tecnici nell'integrazione di PIT nel progetto. Il primo (errore
\texttt{MINION\_DIED}) era dovuto al placeholder \texttt{\$\{argLine\}}
di Surefire, valorizzato solo quando il goal \texttt{jacoco:prepare-agent}
viene eseguito nella stessa invocazione Maven: lanciando PIT da solo,
il placeholder restava non risolto, causando un errore nei processi
di misurazione di PIT. È stata dichiarata una property di progetto
\texttt{<argLine></argLine>} vuota per risolvere il problema. Il
secondo problema riguardava una mutation coverage 0\%/0\% falsa su
\texttt{Utf8}: PIT esegue di base solo test JUnit 3/4, mentre
\texttt{TestUtf8CategoryPartition} è scritta in JUnit 5. È stato
necessario aggiungere la dipendenza \texttt{pitest-junit5-plugin} al
plugin.

\paragraph{Risultati iniziali} La \tabref{tab:pit-results} riassume i risultati numerici, mentre la
\tabref{fig:pit-report} mostra il report generato da PIT. Sulla suite manuale di \texttt{Resolver}, PIT
non ha rilevato alcun mutante sopravvissuto tra quelli coperti
(Test Strength 100\%). Su \texttt{Utf8} sono invece emersi 12 mutanti
sopravvissuti, concentrati sui metodi \texttt{equals()},
\texttt{hashCode()}, \texttt{toString()} e \texttt{setByteLength()}.

\paragraph{Test di rinforzo (MT)} Sono stati scritti 14 nuovi test
manuali mirati a uccidere i 12 mutanti sopravvissuti su \texttt{Utf8}.
Durante la loro progettazione sono emerse due scoperte metodologiche:
un primo tentativo di verifica tramite confronto per identità di
riferimento (\texttt{assertSame}) su una stringa vuota non uccideva il
mutante corrispondente, poiché la JDK ottimizza tale caso restituendo
comunque il letterale internato indipendentemente dal ramo di codice
eseguito -- risolto verificando lo stato interno tramite reflection;
inoltre, un test già presente nella suite BB non esercitava
effettivamente il metodo \texttt{equals()} di \texttt{Utf8}, poiché
il contratto di JUnit 5 poteva invocare \texttt{String.equals()}
anziché \texttt{Utf8.equals()} per l'argomento fornito.

Dei 12 mutanti target, 7 sono stati uccisi dai nuovi test. I restanti
5 sono stati argomentati come mutanti \emph{equivalenti}: i due
percorsi di calcolo di \texttt{equals()} e \texttt{hashCode()} (un
ramo ottimizzato tramite \texttt{Arrays.equals}/\texttt{Arrays.hashCode}
per array esattamente dimensionati, e un loop manuale byte-per-byte
per gli altri casi) calcolano matematicamente lo stesso risultato per
qualunque input valido, verificato empiricamente con test mirati
proprio al confine tra i due rami. Dopo l'integrazione dei test MT, la
mutation coverage di \texttt{Utf8} è salita al 56\% (33/59), con Test
Strength dell'87\% (33/38).

\subsection{Stima della reliability}

La reliability di un sistema può essere stimata come
$reliability = 1 - PFD$, dove $PFD$ (\emph{probability of failure on
demand}) rappresenta la probabilità che un'operazione fallisca,
assumendo un \emph{profilo operazionale} che assegna una probabilità
di utilizzo a ciascuna operazione del sistema. In questo lavoro è
stato adottato un profilo operazionale a probabilità uniforme.

La scelta di quali operazioni considerare è stata guidata da un
criterio di rappresentatività: anziché trattare ogni singolo test
generato automaticamente come un'operazione a sé -- scelta che
sovra-rappresenterebbe input casuali o degeneri rispetto a scenari
d'uso plausibili, specialmente per i test prodotti da Randoop -- sono
state utilizzate le categorie della suite di test manuali (BB) come
operazioni rappresentative, ciascuna pesata uniformemente $1/N$.

Su \texttt{Resolver}, le 8 categorie individuate risultano tutte
associate a test passanti, con una Test Strength del 100\% rilevata
dal mutation testing: la reliability stimata è pertanto pari a 1.0.
Analogamente, su \texttt{Utf8} le 6 categorie risultano tutte associate
a test passanti (suite BB e MT complessivamente), con una reliability
stimata anch'essa pari a 1.0.

Questo risultato va qualificato: rappresenta la reliability osservata
rispetto al profilo operazionale definito dalle categorie testate, non
una garanzia di correttezza sull'intero dominio di input della classe.
L'analisi di mutation testing (Sezione~\ref{sec:mutation}) ha già
mostrato che una parte dei mutanti generati da PIT non è coperta da
alcun test della suite manuale: per le porzioni di codice
corrispondenti non esiste evidenza sperimentale, e la stima di
reliability vale quindi solo entro i limiti del profilo operazionale
adottato.